% ============================================================
% STRUCTURAL IDENTIFIERS
% Book I: Ontology of Viability
% ============================================================

\cleardoublepage

\section*{Structural Identifiers}

\phantomsection
\addcontentsline{toc}{section}{Structural Identifiers}

% ------------------------------------------------------------
% 1. PURPOSE AND SCOPE
% ------------------------------------------------------------

\subsection*{1. Purpose and Scope}

Structural identifiers provide a stable system for
identifying the hierarchical locations of textual
elements within the Viability Philosophy System.

Their primary function is to make the organization
of a book explicitly traceable. They identify where
a particular textual element belongs within the
structural hierarchy, independently of its page
number or physical position in a particular
compiled edition.

A structural identifier refers to a location in
the organization of the text. It does not, by itself,
identify the conceptual or argumentative entity
discussed at that location.

The distinction is fundamental. A chapter, section,
subsection, or paragraph has a structural location.
An axiom, definition, proposition, theory, hypothesis,
argument, or philosophical framework has a logical
identity. These two forms of identification serve
different purposes and must remain distinct.

Structural identifiers are intended to support:

\begin{itemize}
    \item precise identification of textual locations;
    \item navigation across the hierarchical structure
    of the book;
    \item stable internal references;
    \item identification of passages independently
    of page numbering;
    \item traceability during manuscript revision;
    \item coordination between structural references
    and logical entity references.
\end{itemize}

This document establishes the format, components,
assignment rules, and maintenance principles of
structural identifiers.

% ------------------------------------------------------------
% 2. GENERAL STRUCTURE OF AN IDENTIFIER
% ------------------------------------------------------------

\subsection*{2. General Structure of an Identifier}

A structural identifier is composed of a sequence
of hierarchical components separated by periods.

The general format is:

\begin{center}
    \texttt{VS.B01.P01.C01.S01.SS01.P001}
\end{center}

Each component identifies a particular level
within the structure of the work.

The components are interpreted in the following
order:

\begin{enumerate}
    \item System identifier;
    \item Book identifier;
    \item Part identifier;
    \item Chapter identifier;
    \item Section identifier;
    \item Subsection identifier;
    \item Paragraph identifier.
\end{enumerate}

The identifier is hierarchical. Each successive
component narrows the structural location within
the preceding component.

A complete identifier therefore represents
a specific location within the organization
of the book rather than an arbitrary sequence
of labels.

Not every structural element requires every
component. An identifier should contain only
the components necessary to identify the
relevant structural level and its position
within the hierarchy.

% ------------------------------------------------------------
% 3. COMPONENTS OF THE IDENTIFIER
% ------------------------------------------------------------

\subsection*{3. Components of the Identifier}

\subsubsection*{3.1. System Identifier: VS}

The prefix \texttt{VS} identifies the
\textit{Viability Philosophy System}.

It establishes the system to which the
corresponding book and its structural
elements belong.

This prefix is retained across the books
of the system.

\subsubsection*{3.2. Book Identifier: B01}

The component \texttt{B01} identifies
the first book of the Viability Philosophy
System.

For the present volume:

\begin{center}
    \texttt{VS.B01}
\end{center}

identifies Book I, \textit{Ontology of
Viability}.

The book component distinguishes the
structural domain of one book from that
of another within the system.

\subsubsection*{3.3. Part Identifier: P01}

The component \texttt{P01} identifies
the first major part of the book.

Part identifiers follow the format:

\begin{center}
    \texttt{P01}, \texttt{P02},
    \texttt{P03}, \texttt{P04}
\end{center}

The identifier corresponds to the
position of the part within the book.

For example:

\begin{center}
    \texttt{VS.B01.P01}
\end{center}

identifies the first part of Book I.

The abbreviation \texttt{P} is used
for the part component. Its meaning
is determined by its position in the
hierarchical identifier.

\subsubsection*{3.4. Chapter Identifier: C01}

The component \texttt{C01} identifies
a chapter within a part.

Chapter identifiers follow the format:

\begin{center}
    \texttt{C01}, \texttt{C02},
    \texttt{C03}, \ldots
\end{center}

A chapter identifier is interpreted
in relation to the part in which
the chapter occurs.

For example:

\begin{center}
    \texttt{VS.B01.P01.C01}
\end{center}

identifies the first chapter of the
first part of Book I.

\subsubsection*{3.5. Section Identifier: S01}

The component \texttt{S01} identifies
a section within a chapter.

Section identifiers follow the format:

\begin{center}
    \texttt{S01}, \texttt{S02},
    \texttt{S03}, \ldots
\end{center}

A section identifier is interpreted
in relation to its containing chapter.

For example:

\begin{center}
    \texttt{VS.B01.P01.C01.S01}
\end{center}

identifies the first section of the
first chapter of the first part.

\subsubsection*{3.6. Subsection Identifier: SS01}

The component \texttt{SS01} identifies
a subsection within a section.

Subsection identifiers follow the format:

\begin{center}
    \texttt{SS01}, \texttt{SS02},
    \texttt{SS03}, \ldots
\end{center}

A subsection identifier is interpreted
in relation to its containing section.

For example:

\begin{center}
    \texttt{VS.B01.P01.C01.S01.SS01}
\end{center}

identifies the first subsection of the
first section of the first chapter of
the first part.

\subsubsection*{3.7. Paragraph Identifier: P001}

The component \texttt{P001} identifies
a paragraph within its containing
subsection.

Paragraph identifiers follow the format:

\begin{center}
    \texttt{P001}, \texttt{P002},
    \texttt{P003}, \ldots
\end{center}

For example:

\begin{center}
    \texttt{VS.B01.P01.C01.S01.SS01.P001}
\end{center}

identifies the first paragraph of
the first subsection of the first
section of the first chapter of
the first part of Book I.

The abbreviation \texttt{P} is used
both for the part component and for
the paragraph component. This does
not create an ambiguity because
their positions and numerical formats
within the hierarchical identifier
are different.

The part component uses two digits,
whereas the paragraph component
uses three digits.

% ------------------------------------------------------------
% 4. HIERARCHICAL INTERPRETATION
% ------------------------------------------------------------

\subsection*{4. Hierarchical Interpretation}

A structural identifier must be read
from left to right, beginning with
the system and book components and
proceeding toward increasingly specific
textual locations.

The following examples illustrate
the successive levels:

\begin{center}
\begin{tabular}{ll}
    \toprule
    \textbf{Identifier} & \textbf{Structural Level} \\
    \midrule
    \texttt{VS.B01} & Book \\
    \texttt{VS.B01.P01} & Part \\
    \texttt{VS.B01.P01.C01} & Chapter \\
    \texttt{VS.B01.P01.C01.S01} & Section \\
    \texttt{VS.B01.P01.C01.S01.SS01}
        & Subsection \\
    \texttt{VS.B01.P01.C01.S01.SS01.P001}
        & Paragraph \\
    \bottomrule
\end{tabular}
\end{center}

Each additional component specifies
a location within the preceding
structural level.

A shorter identifier designates
a broader structural unit, while
a longer identifier designates
a more specific location.

The hierarchy must remain consistent.
A subsection cannot be identified
without its containing section,
and a paragraph cannot be identified
without the structural components
that locate it within the book.

% ------------------------------------------------------------
% 5. NUMBERING AND ASSIGNMENT RULES
% ------------------------------------------------------------

\subsection*{5. Numbering and Assignment Rules}

Structural identifiers are assigned
according to the actual organization
of the manuscript.

The following rules govern their
assignment.

\subsubsection*{5.1. Sequential Numbering}

Components at the same structural
level are numbered sequentially
within their respective parent unit.

Chapters are numbered within their
parts, sections within their chapters,
subsections within their sections,
and paragraphs within their
subsections.

The numbering of a child component
is therefore interpreted within
the scope of its immediate parent.

\subsubsection*{5.2. Fixed Digit Width}

Part, chapter, section, and subsection
components use two-digit numbering.

Paragraph components use three-digit
numbering.

Leading zeros must be retained.

For example:

\begin{itemize}
    \item \texttt{P01}, not \texttt{P1};
    \item \texttt{C03}, not \texttt{C3};
    \item \texttt{S07}, not \texttt{S7};
    \item \texttt{SS02}, not \texttt{SS2};
    \item \texttt{P009}, not \texttt{P9}.
\end{itemize}

This convention preserves a consistent
identifier format and facilitates
systematic sorting and retrieval.

\subsubsection*{5.3. Local Sequential Numbering}

Numbering at each level begins
within its own parent unit.

For example, the first section
of a new chapter may be identified
as \texttt{S01}, even if the preceding
chapter also contains a section
identified as \texttt{S01}.

The complete structural identifier
distinguishes the two locations
through their different parent
components.

\subsubsection*{5.4. Assignment to Paragraphs}

Paragraph identifiers are assigned
to paragraphs that require precise
structural reference.

The identifier must correspond
to the paragraph's actual position
within its containing subsection.

Paragraph identification should
not be confused with sentence
numbering, line numbering, or
page numbering.

A paragraph is a textual structural
unit, not an individual sentence
or typographical line.

% ------------------------------------------------------------
% 6. STRUCTURAL IDENTIFIERS AND PAGE NUMBERS
% ------------------------------------------------------------

\subsection*{6. Structural Identifiers and Page Numbers}

Structural identifiers are independent
of page numbering.

Page numbers may change as a result
of editing, typesetting, pagination,
or changes in the physical edition
of the book.

Structural identifiers, by contrast,
refer to the hierarchical organization
of the manuscript.

A reference to a structural identifier
should therefore remain meaningful
even when the corresponding passage
moves to a different page.

Page numbers may be supplied as
additional navigational information,
but they must not replace structural
identifiers where precise structural
identification is required.

The two systems serve complementary
purposes:

\begin{itemize}
    \item structural identifiers
    identify textual locations
    within the hierarchy;

    \item page numbers identify
    physical locations within
    a particular edition.
\end{itemize}

% ------------------------------------------------------------
% 7. STRUCTURAL IDENTIFIERS AND LOGICAL
% ENTITY IDENTIFIERS
% ------------------------------------------------------------

\subsection*{7. Structural Identifiers and
Logical Entity Identifiers}

Structural identifiers and logical
entity identifiers must remain
distinct throughout the book.

A structural identifier identifies
the location at which a conceptual
or argumentative entity is presented.

A logical entity identifier identifies
the entity itself, independently
of any single textual location.

For example:

\begin{center}
\begin{tabular}{ll}
    \toprule
    \textbf{Identifier} & \textbf{Identification Type} \\
    \midrule
    \texttt{VS.B01.P01.C01.S01.SS01.P001}
        & Structural location \\
    \texttt{AX-001}
        & Logical entity \\
    \bottomrule
\end{tabular}
\end{center}

The first identifier locates a paragraph.
The second identifies a logical entity.

The entity \texttt{AX-001} may be
introduced in one paragraph, explained
in another, and used in an argument
elsewhere. Its logical identifier
remains the same across these locations.

The structural identifiers of those
locations remain distinct.

The rules governing logical entity
identification are established in
the document \textit{Logical Entity
Identifiers}.

% ------------------------------------------------------------
% 8. STABILITY AND REVISION
% ------------------------------------------------------------

\subsection*{8. Stability and Revision}

Structural identifiers must be
maintained consistently throughout
the development of the manuscript.

However, structural identifiers
are tied to the actual hierarchy
of the text. A change in that
hierarchy may require a corresponding
change in the identifier.

The following principles govern
such changes.

\subsubsection*{8.1. Changes in Wording}

A revision of a title, sentence,
or paragraph that does not alter
the structural position of the
element does not, by itself,
require a change in its structural
identifier.

\subsubsection*{8.2. Relocation of Structural Elements}

When a chapter, section, subsection,
or paragraph is moved to another
position in the hierarchy, its
structural identifier must be
reassessed according to the new
location.

The identifier must accurately
represent the element's current
structural position.

\subsubsection*{8.3. Insertion of New Elements}

When a new structural element is
inserted, it must receive an
identifier consistent with the
numbering rules of its level.

Existing identifiers should be
preserved wherever possible.
However, if the insertion changes
the sequential structure, the
affected identifiers must be
reviewed to maintain consistency.

\subsubsection*{8.4. Removal of Elements}

When a structural element is
removed, its identifier should
not be reassigned casually to
an unrelated element.

References to the removed location
must be reviewed and corrected
where necessary.

\subsubsection*{8.5. Revision of Referenced Locations}

Whenever a structural identifier
changes, all internal references
that depend on that identifier
must be checked.

No reference should remain directed
toward an obsolete structural
location when the intended element
has been moved, renamed, or removed.

% ------------------------------------------------------------
% 9. STRUCTURAL IDENTIFIERS AND INTERNAL
% DOCUMENT ORGANIZATION
% ------------------------------------------------------------

\subsection*{9. Structural Identifiers and
Internal Document Organization}

The structural hierarchy of the
book is represented through its
parts, chapters, sections,
subsections, and paragraphs.

The identifier system follows
this hierarchy rather than
the names of source files used
to compose the manuscript.

A source file may contain one
or more structural elements,
and a structural element may
be revised or relocated without
requiring that its identity be
defined by the name of a source
file.

File names and directory paths
are organizational features of
the manuscript project. They
must not be confused with
structural identifiers.

The structural identifier
therefore remains a reference
to the textual architecture
of the book, not to a particular
file-system location.

% ------------------------------------------------------------
% 10. EXTERNAL SOURCES
% ------------------------------------------------------------

\subsection*{10. External Sources}

External books, articles, documents,
and other sources do not receive
structural identifiers belonging
to the internal hierarchy of
the Viability Philosophy System.

They retain their own bibliographic
identities and are referenced
according to the conventions
established for external sources.

The presence of an external source
within a particular paragraph
does not make that source a
structural component of the book.

Similarly, a quotation or citation
does not acquire a structural
identifier independently of
the textual location in which
it appears.

The distinction between internal
structural identification and
external bibliographic identification
must be maintained throughout
the work.

% ------------------------------------------------------------
% 11. SUMMARY OF THE STRUCTURAL
% IDENTIFIER FORMAT
% ------------------------------------------------------------

\subsection*{11. Summary of the Structural
Identifier Format}

The structural identifier format
for Book I is illustrated by:

\begin{center}
    \texttt{VS.B01.P01.C01.S01.SS01.P001}
\end{center}

Its components are:

\begin{center}
\begin{tabular}{lll}
    \toprule
    \textbf{Component} & \textbf{Meaning}
        & \textbf{Example} \\
    \midrule
    VS & System & \texttt{VS} \\
    B01 & Book & \texttt{B01} \\
    P01 & Part & \texttt{P01} \\
    C01 & Chapter & \texttt{C01} \\
    S01 & Section & \texttt{S01} \\
    SS01 & Subsection & \texttt{SS01} \\
    P001 & Paragraph & \texttt{P001} \\
    \bottomrule
\end{tabular}
\end{center}

The identifier must be interpreted
hierarchically and assigned
according to the actual structure
of the manuscript.

Its purpose is to make textual
locations precise, stable enough
for systematic referencing, and
independent of changes in page
numbering.

Structural identification
constitutes one component of
the broader referencing system.
It must be used in coordination
with logical entity identifiers,
internal cross-references,
entity location references,
logical dependencies, and
bibliographic references.

\cleardoublepage

% ============================================================
% END OF STRUCTURAL IDENTIFIERS
% ============================================================